\subsection{Function basics}
\tsDef{Function}{
    A function (mapping) is a rule that assigns to each admissible input exactly one output in the target set.
    The set of all possible inputs is called the \emph{domain}, denoted by $\mathrm{domain}(f)$ or $D$.
    The set of all actually attained outputs is called the \emph{image}, denoted by $\mathrm{image}(f)$ or $W$.
    \[
        f : \mathrm{domain}(f) \to \mathrm{range}(f),
        \quad
        x \mapsto f(x).
    \]
    The input is called the \emph{independent variable}, the output the \emph{dependent variable}.
}
\tsDef{Restriction of a Function}{
    Let $f : D(f) \to \mathbb{R}$ and let $D_0 \subseteq D(f)$.
    The restriction of $f$ to $D_0$ is the function
    \[
        f|_{D_0} : D_0 \to \mathbb{R},
        \quad
        f|_{D_0}(x) = f(x).
    \]
    Note that $f$ and $f|_{D_0}$ are, in general, different functions.
}
\tsDef{Bounded Functions}{
    Let $f : D(f) \to \mathbb{R}$.
    \begin{itemize}[noitemsep]
        \item $f$ is \emph{bounded above} if $\exists M \in \mathbb{R}$ such that
              \(
                  f(x) \le M \quad \forall x \in D(f).
              \)
        \item $f$ is \emph{bounded below} if $\exists M \in \mathbb{R}$ such that
              \(
                  f(x) \ge M \quad \forall x \in D(f).
              \)
        \item $f$ is \emph{bounded} if $\exists M \in \mathbb{R}$ such that
              \(
                  |f(x)| \le M \quad \forall x \in D(f).
              \)
    \end{itemize}
}
\tsDef{Injective, Surjective, Bijective}{
    Let $f : X \to Y$.
    \begin{itemize}[noitemsep]
        \item $f$ is \emph{injective} if
              \(
                  f(x_1) = f(x_2) \Rightarrow x_1 = x_2.
              \)
        \item $f$ is \emph{surjective} if
              \(
                  \forall y \in Y \; \exists x \in X : f(x) = y.
              \)
        \item $f$ is \emph{bijective} if it is both injective and surjective.
    \end{itemize}
}
\tsDef{Composition}{
    Let $f : X \to Y$ and $g : Y \to Z$.
    The composition is defined by
    \[
        g \circ f : X \to Z,
        \quad
        (g \circ f)(x) = g(f(x)).
    \]
}
\tsDef{Identity and Inverse}{
    For a set $X$, the identity function is defined by
    \(
        \mathrm{id}_X(x) = x \ (\forall x \in X).
    \)
    If $f : X \to Y$ is bijective, then there exists a unique function
    \(
        f^{-1} : Y \to X
    \)
    such that
    \(
        f^{-1} \circ f = \mathrm{id}_X,
        \
        f \circ f^{-1} = \mathrm{id}_Y.
    \)
    This function is called the \emph{inverse} of $f$.
}
\tsDef{Monotonicity}{
    Let $f : D(f) \subseteq \mathbb{R} \to \mathbb{R}$.
    \begin{itemize}[noitemsep]
        \item $f$ is \emph{monotonically increasing} if
              \(
                  x_1 < x_2 \Rightarrow f(x_1) \le f(x_2).
              \)
        \item $f$ is \emph{strictly increasing} if
              \(
                  x_1 < x_2 \Rightarrow f(x_1) < f(x_2).
              \)
        \item $f$ is \emph{monotonically decreasing} if
              \(
                  x_1 < x_2 \Rightarrow f(x_1) \ge f(x_2).
              \)
        \item $f$ is \emph{strictly decreasing} if
              \(
                  x_1 < x_2 \Rightarrow f(x_1) > f(x_2).
              \)
    \end{itemize}
}
\tsThe{Strict Monotonicity Implies Injectivity}{
    Every strictly monotone function is injective.
}
\tsDef{Even and Odd Functions}{
    Let $f : \mathbb{R} \to \mathbb{R}$.
    \begin{itemize}[noitemsep]
        \item $f$ is \emph{even} if
              \(
                  f(-x) = f(x).
              \)
        \item $f$ is \emph{odd} if
              \(
                  f(-x) = -f(x).
              \)
    \end{itemize}
}
\tsDef{Convexity and Concavity (Geometric)}{
    The graph of a function is called
    \begin{itemize}[noitemsep]
        \item \emph{convex} if it bends to the left when moving from left to right,
        \item \emph{concave} if it bends to the right when moving from left to right.
    \end{itemize}
}
\subsection{Limits of Functions}
\tsDef{Limit of a Function ($\varepsilon$-$\delta$ definition)}{
    Let $f : D(f) \to \mathbb{R}$ and $x_0 \in \mathbb{R}$.
    A number $L \in \mathbb{R}$ is called the limit of $f(x)$ at $x_0$ if
    \[
        \forall \varepsilon > 0 \; \exists \delta > 0 \; \text{such that}
        \quad
        x \in D(f), \; |x - x_0| < \delta
        \Rightarrow
        |f(x) - L| < \varepsilon.
    \]
    The limit is uniquely determined.
}
\tsDef{One-Sided Limits}{
    One-Sided Limits
    \begin{itemize}[noitemsep]
        \item Right-hand limit:
              \(
                  \lim_{x \to x_0^+} f(x) = L
              \)
        \item Left-hand limit:
              \(
                  \lim_{x \to x_0^-} f(x) = L
              \)
    \end{itemize}
}
\tsDef{Limits at Infinity}{
    Limits at Infinity
    \begin{itemize}[noitemsep]
        \item
              \(
              \lim_{x \to \infty} f(x) = L
              \)
        \item
              \(
              \lim_{x \to -\infty} f(x) = L
              \)
    \end{itemize}
}
\tsDef{Infinite Limits}{
    Infinite Limits
    \begin{itemize}[noitemsep]
        \item
              \(
              \lim_{x \to c} f(x) = \infty
              \)
        \item
              \(
              \lim_{x \to c} f(x) = -\infty
              \)
    \end{itemize}
}
\tsThe{Limit Laws (Algebra of Limits for Functions)}{
    Let
    \[
        \lim_{x \to c} f(x) = K,
        \quad
        \lim_{x \to c} g(x) = L,
        \quad
        K,L \neq \pm\infty.
    \]
    Then:
    \begin{itemize}[noitemsep]
        \item
              \(
              \lim_{x \to c} (f+g) = K+L
              \)
        \item
              \(
              \lim_{x \to c} (f-g) = K-L
              \)
        \item
              \(
              \lim_{x \to c} (fg) = KL
              \)
        \item
              \(
              \lim_{x \to c} (A f) = A K
              \)
        \item If $L \neq 0$:
              \(
              \lim_{x \to c} \frac{f}{g} = \frac{K}{L}
              \)
    \end{itemize}
}
\columnbreak
\subsection{Continuity of functions}
\tsDef{Continuity ($\varepsilon$-$\delta$ definition)}{
    Let $f : D \to \mathbb{R}$ and $x_0 \in D$. The function $f$ is called \emph{continuous at $x_0$} if
    \[
        \forall \varepsilon > 0 \; \exists \delta > 0 \;\text{such that}\;
        |x - x_0| < \delta \;\Rightarrow\; |f(x) - f(x_0)| < \varepsilon.
    \]
    The function is called \emph{continuous} if it is continuous at every point in its domain.
}
\tsDef{One-Sided Continuity}{
    Let $f : D(f) \to \mathbb{R}$ and $x_0 \in D(f)$.
    \begin{itemize}[noitemsep]
        \item $f$ is \emph{right-continuous} at $x_0$ if
              \(
                  \lim_{x \to x_0^+} f(x) = f(x_0).
              \)
        \item $f$ is \emph{left-continuous} at $x_0$ if
              \(
                  \lim_{x \to x_0^-} f(x) = f(x_0).
              \)
    \end{itemize}
}
\tsThe{Sequential Characterization of Continuity (Heine's Theorem)}{
    Let $f : D \subseteq \mathbb{R} \to \mathbb{R}$ and $x_0 \in D$.
    Then $f$ is continuous at $x_0$ if and only if for every sequence $(x_n)$ with $x_n \to x_0$ we have
    \(
        f(x_n) \to f(x_0).
    \)
}
\tsThe{Algebra of Continuous Functions}{
    Let $f,g : D \subseteq \mathbb{R} \to \mathbb{R}$ be continuous. Then:
    \begin{itemize}[noitemsep]
        \item $f+g$ is continuous
        \item $f-g$ is continuous
        \item $c f$ is continuous for any constant $c \in \mathbb{R}$
        \item $f \cdot g$ is continuous
        \item $\frac{f}{g}$ is continuous if $g \neq 0$
    \end{itemize}
}
\tsThe{Continuity of Composition}{
    Let $f : I \to \mathrm{image}(f)$ and $g : \mathrm{image}(f) \to \mathbb{R}$ be continuous.
    Then the composition
    \(
        g \circ f
    \)
    is continuous and
    \[
        \lim_{x \to x_0} g(f(x)) = g\!\left( \lim_{x \to x_0} f(x) \right).
    \]
}
\tsThe{Continuity of the Inverse Function}{
    Let $I$ be an interval and let $f : I \to J$ be continuous and bijective.
    Then $J$ is also an interval and the inverse function
    \(
        f^{-1} : J \to I
    \)
    is continuous.
}
\tsThe{Intermediate Value Theorem}{
Let $f : [a,b] \to \mathbb{R}$ be continuous and let $c$ be a value between $f(a)$ and $f(b)$.
Then there exists $x \in [a,b]$ such that
\(
    f(x) = c.
\)
}
\tsCor{Zero of a Continuous Function (Bolzano theorem)}{
    Let $f : I \subseteq \mathbb{R} \to \mathbb{R}$ be continuous and suppose
    \[
        f(a) < 0, \quad f(b) > 0.
    \]
    Then there exists $x \in (a,b)$ such that
    \[
        f(x) = 0.
    \]
}
\tsDef{Compact Interval}{
    A bounded and closed interval $[a,b]$ is called \emph{compact}.
}
\tsLem{Bolzano--Weierstrass Theorem (for Compact Intervals)}{
    Let $(x_n)$ be a sequence in a compact interval $[a,b]$.
    Then there exists a subsequence $(x_{n_k})$ such that
    \(
        x_{n_k} \to x_0
        \quad \text{with } x_0 \in [a,b].
    \)
}
\tsDef{Local Extrema}{
    Let $f : D \subseteq \mathbb{R} \to \mathbb{R}$.
    \begin{itemize}[noitemsep]
        \item $x_0$ is a \emph{local maximum} if
              \(
                  f(x_0) \ge f(x) \quad \forall x \in D.
              \)
        \item $x_0$ is a \emph{local minimum} if
              \(
                  f(x_0) \le f(x) \quad \forall x \in D.
              \)
        \item Local maxima and minima are called \emph{local extrema}.
    \end{itemize}
}
\tsThe{Extreme Value Theorem (Weierstrass)}{
    Let $f : I \subseteq \mathbb{R} \to \mathbb{R}$ be continuous and let $I$ be compact.
    Then $f$ is bounded and attains its maximum and minimum on $I$, i.e.,
    there exist $x_{\min}, x_{\max} \in I$ such that
    \(
        f(x_{\min}) \le f(x) \le f(x_{\max})
        \ (\forall x \in I).
    \)
}